%&latex
\documentclass[a4paper]{article}

\frenchspacing

\usepackage[cp1250]{inputenc}
\usepackage[czech]{babel}

\usepackage{a4wide}
\usepackage{amsmath, amsthm, amssymb, amsfonts}
\usepackage[mathcal]{eucal}
\usepackage{graphicx}
\usepackage{url}
\usepackage{color}
\usepackage{wrapfig}
\usepackage{capt-of}
\usepackage{float}
\usepackage{ifthen}
\usepackage{subfig}
\usepackage[normalem]{ulem}


% sirka a vyska textu nastavena jako papir, vsechny okraje vynulovany a pridano 20pt na kazdou stranu
% horizontalni rozmery
\setlength{\textwidth}{\paperwidth}
\addtolength{\textwidth}{-40pt}
\addtolength{\hoffset}{-1in}
\addtolength{\hoffset}{20pt}
\setlength{\oddsidemargin}{0in}
\setlength{\marginparsep}{0in}
% vertikalni rozmery
\setlength{\textheight}{\paperheight}
\addtolength{\textheight}{-60pt}
\addtolength{\voffset}{-1in}
\addtolength{\voffset}{20pt}
\setlength{\topmargin}{0in}
\setlength{\headheight}{0in}
\setlength{\headsep}{0in}


%Obrazek na miste
%pouziti
%%\obrazeknahore{adresa}{popisek}{label}
\long\def\obrazeknahore#1#2#3 {

\begin{figure}[t]
    \centering
    \includegraphics[width=0.8\textwidth]{#1}
    
    \caption{#2}
    \label{#3}
    
\end{figure}

}


%==========================================
%PEKELNA MAKRA NA ZAROVNANI OBRAZKU DOPRAVA

\makeatletter


%tohle je makro, ktere mi dovoluje obtekani i u kratkych environmentu
%ABSOLUTNE nechapu, jak to funguje, ale funguje to
%viz http://tex.stackexchange.com/questions/26078/ 
\def\odrovnej{\@@par
\ifnum\@@parshape=\z@ \let\WF@pspars\@empty \fi % reset `parshape'
\global\advance\c@WF@wrappedlines-\prevgraf \prevgraf\z@
\ifnum\c@WF@wrappedlines<\tw@ \WF@finale \fi}

\makeatother



%---
%makro, co da obrazek doprava a ostatni text ho obteka
%(bez toho predchazejiciho makra to ale poradne nebeha)
%pouziti:
%\obrazekvpravo{adresa}{popisek}{label}{procento sirky}
\long\def\obrazekvpravo#1#2#3#4{

\setlength\intextsep{-20pt}

    \begin{wrapfigure}{r}{#4\textwidth}
      \begin{center}
          \vspace{-10pt}
          
        \includegraphics[width=#4\textwidth]{#1}
        \vspace{-10pt}
        
      \end{center}
      
      \caption{#2}
      \label{#3}
      
      
    \end{wrapfigure}

\setlength\intextsep{0pt}

    
}




%---
%makro pro pripady, kdy wrapfigure neco mrsi
%je to docela pekelne
%je nutne mu dat jak text vpravo, tak text vlevo
%a nevim, jestli bude 100% fungovat, ale doufam, ze jo

%pouziti:
%\obrazekvpravominipage{adresa}{popisek}{label}{procento sirky}{1 - procento sirky}{text vlevo}
\long\def\obrazekvpravominipage#1#2#3#4#5#6{

\noindent\begin{minipage}{#5\linewidth}
\vspace{0pt}
#6
\end{minipage}
\hspace{0.5cm}
\noindent\begin{minipage}{#4\linewidth}
\vspace{0pt}
\centering
\includegraphics[width=0.9\textwidth]{#1}
\captionof{figure}{#2}
\label{#3}
\end{minipage}

}

%KONEC PEKELNYCH MAKER
%=====================

\def\lebka{\includegraphics[width=3mm]{../common/lebka}}

% makra pro poznamku u vyrokove a predikatove logiky
\def\vl{ -- ve v�rokov� logice}
\def\pl{ -- v predik�tov� logice}


%Vacsina prostredi je dvojjazicne. V pripade, ze znenie napr pozorovania je pisane po slovensky, malo by byt po slovensky aj oznacenie.

\def\ifis#1#2{\ifthenelse{\equal{#1}{0}}{}{#2}}

%\newenvironment{e}[3]{\pagebreak[2]\noindent\textbf{\ifis{#2}{$\bigstar$}#1}\ifis{#3}{\emph{~(#3)}}\par\noindent\leftskip 10pt }{\odrovnej\par\bigskip}
\newenvironment{e}[3]{\pagebreak[2]\noindent\textbf{\ifis{#2}{\lebka~}#1}\ifis{#3}{\emph{~(#3)}}\ifis{#2}{~\lebka}\par\noindent\leftskip 10pt }{\odrovnej\par\bigskip}


%obecne prostredie, ktore ma vyuzitie pri specialnych odstavcoch ako (uloha, algoritmus...) aby nevzniklo dalsich x prostredi
\newenvironment{obecne}[1]{\pagebreak[2]\noindent\textbf{#1}\par\noindent\leftskip 10pt}{\odrovnej\par\bigskip}

\newenvironment{report}{\pagebreak[2]\noindent\textbf{Report}\em\par\noindent\leftskip 10pt}{\par\bigskip}

%\newenvironment{reportN}[1]{\pagebreak[2]\noindent\textbf{Report~}\emph{(#1)}\emph\par\noindent\leftskip 10pt}{\odrovnej\par\bigskip}
\newenvironment{reportN}[1]{\pagebreak[2]\noindent\textbf{Report~}\emph{(#1)}\em\par\noindent\leftskip 10pt}{\odrovnej\par\bigskip}

\newenvironment{penumerate}{
\begin{enumerate}
  \setlength{\itemsep}{1pt}
  \setlength{\parskip}{0pt}
  \setlength{\parsep}{0pt}
  %\setlength{\topsep}{200pt}
  \setlength{\partopsep}{200pt}
}{\end{enumerate}}

\def\pismenka{\numberedlistdepth=2} %pouzit, ked clovek chce opismenkovany zoznam...

\newenvironment{pitemize}{
\begin{itemize}
  \setlength{\itemsep}{1pt}
  \setlength{\parskip}{0pt}
  \setlength{\parsep}{0pt}
}{\end{itemize}}

\definecolor{gris}{gray}{0.95}
\newcommand{\ramcek}[2]{\begin{center}\fcolorbox{white}{gris}{\parbox{#1}{#2}}\end{center}\par}


\title{\LARGE U�ebn� texty k st�tn� bakal��sk� zkou�ce \\ Matematika \\ Metrick� prostory}

\begin{document}

\maketitle

\newpage
\setcounter{section}{5}
\section{Metrick� prostory}

\begin{e}{Po�adavky}{0}{0}
\begin{pitemize}
    \item Definice metrick�ho prostoru, p��klady.
    \item Spojitost a stejnom�rn� spojitost.
    \item Kompaktn� prostory a jejich vlastnosti, �pln� prostory.
\end{pitemize}
\end{e}

\subsection{Definice metrick�ho prostoru, p��klady}

\subsubsection*{Metrick� prostor, metrika}

\begin{definiceN}{metrick� prostor}
\emph{Metrick� prostor} je dvojice $(M, d)$, kde $M$ je mno�ina a $d: M \times M \rightarrow \mathbb{R}$ je zobrazen�, zvan� \emph{metrika}, spl�uj�c� t�i axiomy:
\begin{penumerate}
    \item $d(x, y) = 0 \Leftrightarrow x = y$
    \item $d(x, y) = d(y, x)$ \hfill \emph{(symetrie)}
    \item $d(x, y) \le d(x, z) + d(z, y)$ \hfill\emph{(troj�heln�kov� nerovnost)}
\end{penumerate}
Metrick� prostory jsou abstrakc� jevu vzd�lenosti. Z axiom� 1. a 3. vypl�v� nez�pornost hodnot metriky (kter� se ale v�t�inou explicitn� uv�d� jako sou��st prvn�ho axiomu). Prvky metrick�ho prostoru nazveme \emph{body}.
\end{definiceN}

\begin{obecne}{P��klady metrik}
Nech� $M = \mathbb{R}^n$ a $p \ge 1$ je re�ln� ��slo. Na $M$ definujeme metriky\\(kde $x = (x_1, x_2, \dots, x_n)$, $y = (y_1, y_2, \dots , y_n)$)
$$d_p(x, y) = \left(\sum_{i=1}^{n}|x_i - y_i|^p\right)^{1/p}$$

\noindent Potom:\begin{pitemize}
\item Pro $p = 1, n = 1$ dost�v�me metriku $|x - y|$.
\item Pro $p = 2, n \ge 2$ dost�v�me euklidovskou metriku
$$d_2(x, y) = \parallel x - y \parallel = \sqrt{(x_1-y_1)^2+(x_2-y_2)^2+\dots+(x_n-y_n)^2}$$
\item Pro $p = 1, n \ge 2$ dost�v�me tzv. po���ckou metriku a pro $p \rightarrow \infty$ maximovou
metriku
$$d_1(x, y) = \max_{1 \le i \le n} |x_i - y_i|$$
\end{pitemize}
\par\medskip\noindent
Bu� $X$ libovoln� mno�ina. $F(X)$ ozna�me mno�inu v�ech omezen�ch funkc� $f:X\to\mathbb{R}$ a definujme funkci $d$ p�edpisem
$$d(f,g)=\sup_{x\in X}|f(x)-g(x)|$$
Pak $(F(X),d)$ je metrick� prostor (s tzv. supremovou metrikou).
\par\medskip\noindent
Vezmeme $M = \mathcal{C}[a, b]$ (mno�ina re�ln�ch funkc� definovan�ch a spojit�ch an intervalu $[a, b]$) a re�ln� ��slo $p \geq 1$. Podobn� jako v prvn�m p��kladu m�me na $M$ metriky
$$d_p\left(f, g\right) = \left( \int_a^b \lvert f(x) - g(x) \rvert^p \; \mathrm{d}x\right)^{1/p}.$$
Pro $p = 1$ m�me integr�ln� metriku a $p \to \infty$ d�v� maximovou metriku
z druh�ho p��kladu. Vezmeme-li �ir�� t��du funkc� $M = \mathcal{R}[a, b]$ (funkce maj�c� na $[a, b]$ Riemann�v
integr�l), je $d_p(f, g)$ definovan�, ale nespl�uje axiom 1) a nedost�v�me metriku.
Zm�n�me-li nap��klad hodnotu funkce $f \in \mathcal{R}[a, b]$ v jedin�m bod�, dostaneme
odli�nou funkci $f_0 \in \mathcal{R}[a, b]$, ale $d_p(f, f_0) = 0$. (Tato pot� se odstran� tak, �e
m�sto s $\mathcal{R}[a, b]$ se pracuje s $\mathcal{R}[a, b]/\sim$ pro vhodnou relaci ekvivalence $\sim$.)
\par\medskip\noindent
Na souvisl�m grafu $G = (M, E)$ s mno�inou vrchol� $M$ m�me metriku $d(u, v)$ odpov�daj�c� po�tu hran na nejkrat�� cest� v $G$ spojuj�c� $u$ a $v$.
\par\medskip\noindent
Je-li $A$ kone�n� mno�ina (abeceda), m�me na mno�in� $M = A^m$ slov d�lky $m$ nad abecedou $A$ tzv. Hammingovu metriku ($u = a_1 a_2 \ldots a_m, v = b_1 b_2 \ldots b_m$), kde $d(u, v)$ odpov�d� po�tu sou�adnic $i$, pro n� $a_i \neq b_i$. M��� m�ru odli�nosti obou slov -- jak� nejmen�� po�et zm�n v p�smenech sta�� k p�em�n� $u$ ve $v$.
\par\medskip\noindent
�pln� trivi�ln� p��klad metriky dostaneme, kdy� na n�jak� mno�in� $M$ polo��me $d(x,y)=1$ pro $x\neq y$ a $d(x,x)=0$.
\end{obecne}

\begin{definiceN}{euklidovsk� prostor}
\emph{Euklidovsk�m prostorem} rozum�me metrick� prostor $(\mathbb{R}^n, d_2)$, kde $d_2$ je funkce dan� p�edpisem $d_2(x,y)=\sqrt{\sum_{i=1}^n(x_i-y_i)^2}$.
\end{definiceN}

\subsubsection*{Konvergence}

\begin{definiceN}{konvergence posloupnosti bod�}
�ekneme, �e posloupnost bod� $(x_n)_{n\geq 0}$ n�jak�ho metrick�ho prostoru $(M, d)$ \emph{konverguje k bodu $x \in M$} ($x_n\to x$), nebo �e $x=\lim_{n\geq 0} x_n$, jestli�e
$$\forall\varepsilon>0\ \exists n_0:n\geq n_0\;\Rightarrow\; d(x_n,x)<\varepsilon.$$
\end{definiceN}

\begin{poznamkaN}{vlastnosti konvergence}
Nech� je d�n metrick� prostor $(M, d)$ a v n�m posloupnost bod� $(x_n)_{n\geq 0}$. Potom plat�:
\begin{penumerate}
    \item Jestli�e pro bod $y\in M$ plat�, �e $\exists n_0\in\mathbb{N}: \forall n\geq n_0: x_n = y$, pak $x_n\to y$
    \item Nech� $x_n\to y_1$ a z�rove� $x_n\to y_2$. Potom $y_1=y_2$.
    \item Podposloupnost konvergentn� posloupnosti konverguje ke stejn�mu bodu.
\end{penumerate}
\begin{dukaz}
Obdobn� jako pro konvergenci posloupnost� z $\mathbb{R}$ -- pouze m�sto absolutn� hodnoty pou�ijeme obecnou metriku.
\end{dukaz}
\end{poznamkaN}
\newpage

\subsection{Spojitost a stejnom�rn� spojitost}
\subsubsection*{Otev�en� a uzav�en� mno�iny}

\begin{definiceN}{otev�en� a uzav�en� koule}
Bu� $(M,d)$ metrick� prostor, $x \in M, r > 0$, potom
\begin{pitemize}
\item \emph{otev�enou koul�} (\emph{$r$-okol�m}) se st�edem $x$ a polom�rem $r$ nazveme mno�inu
$$B(x,r)=\{y \in M | d(x,y)<r\}$$
\item \emph{uzav�enou koul�} (\emph{$r$-okol�m}) se st�edem $x$ a polom�rem $r$ nazveme mno�inu
$$\overline{B}(x,r) = \{y \in M | d(x,y) \le r \}$$
\end{pitemize}
\end{definiceN}

\begin{definiceN}{otev�en� a uzav�en� mno�ina}
Bu� $(M,d)$ metrick� prostor, $G \subseteq M$, potom
\begin{pitemize}
\item $G$ je \emph{otev�en�} v $M$, pokud
$$\forall x \in G\;\exists r >\;0:\; B(x,r)\subseteq G$$
(tj. mno�ina $G$ je okol�m ka�d�ho sv�ho bodu)
\item $G$ je \emph{uzav�en�} v $M$, pokud jej� dopln�k $M \setminus G$ je otev�en� v $M$.
\end{pitemize}
\par\medskip\noindent
Otev�en� koule je otev�en� mno�ina v ka�d�m metrick�m prostoru. Podobn� pro uzav�enou kouli.
\end{definiceN}

\begin{vetaN}{vlastnosti otev�en�ch mno�in}
Bu� $(M,d)$ metrick� prostor, potom plat�
\begin{penumerate}
\item $\emptyset$, $M$ jsou otev�en� v $M$
\item kone�n� pr�nik otev�en�ch mno�in je otev�en� mno�ina v $M$
\item libovoln� velk� sjednocen� otev�en�ch mno�in je otev�en� mno�ina v $M$
\end{penumerate}
\end{vetaN}

\begin{vetaN}{vlastnosti uzav�en�ch mno�in}
Nech $(M,d)$ je metrick� prostor, potom plat�
\begin{penumerate}
\item $\emptyset$, $M$ jsou uzav�en� v $M$
\item libovoln� pr�nik uzav�en�ch mno�in je uzav�en� mno�ina v $M$
\item kone�n� sjednocen� uzav�en�ch mno�in je uzav�en� mno�ina v $M$
\end{penumerate}
\end{vetaN}
\begin{definiceN}{okol� bodu, vnit�n�, vn�j�� a dal�� body}
Bu� $(M, d)$ metrick� prostor, $a \in M$ jeho bod. Ka�dou otev�enou mno�inu $U \subset M$ spl�uj�c� $a \in U$ nazveme \emph{okol�m bodu $a$}. Tedy ka�d� koule $B(a, r)$ je okol�m bodu $a$ a ov�em ka�d� okol�
bodu $a$ obsahuje n�jakou takovou kouli. V n�sleduj�c�ch definic�ch bu�te $a \in M$ bod, $X \subset M$ mno�ina a $U$ okol� bodu $a$. �ekneme, �e
\begin{pitemize}
    \item $a$ je \emph{vnit�n�m} bodem $X$, kdy� existuje $U$ tak, �e $U \subset X$,
    \item $a$ je \emph{vn�j��m} bodem $X$, kdy� existuje $U$ tak, �e $U \subset M\setminus X$,
    \item $a$ je \emph{hrani�n�m} bodem $X$, kdy� ka�d� $U$ prot�n� $X$ i $M\setminus X$,
    \item $a$ je \emph{limitn�m} bodem $X$, kdy� je pro ka�d� $U$ pr�nik $U \cap X$ nekone�n�,
    \item $a$ je \emph{izolovan�m} bodem $X$, kdy� existuje $U$ tak, �e $U \cap X = \lbrace a\rbrace$.
\end{pitemize}
\end{definiceN}

\begin{priklad}
Vnit�n� a izolovan� body $X$ nutn� le�� v $X$, vn�j�� body le�� mimo $X$. Hrani�n� a limitn� body $X$ mohou le�et v $X$ i mimo $X$. Vezm�me v euklidovsk� rovin� $R^2$ mno�inu $X = \lbrace x \in R^2\;\vert\;0 < \|x\| < 1\rbrace \cup \lbrace(0, 2)\rbrace$, jednotkov� kruh se st�edem v po��tku, z n�ho� jsme po��tek odstranili a k
n�mu� jsme p�idali bod $(0, 2)$. Pak vnit�n� body $X$ tvo�� mno�inu $\lbrace x \in R^2\;\vert\;0 <
\|x\| < 1\rbrace$, vn�j�� body mno�inu $\lbrace x \in R^2\;\vert\; \|x\| > 1, x \neq (0, 2)\rbrace$, hrani�n�
body mno�inu $\lbrace x \in R^2\;\vert\; \|x\| = 1\rbrace \cup \lbrace(0, 0), (0, 2)\rbrace$, limitn� body mno�inu
$\lbrace x \in R^2\;\vert\;\|x\| \leq 1\rbrace$ a izolovan� body mno�inu $\lbrace(0, 2)\rbrace$.
\end{priklad}

\begin{definiceN}{vzd�lenost bodu od mno�iny}
V metrick�m prostoru $(X, d)$ bu� $A\subset X$ mno�ina a $x\in(X, d)$ bod. \emph{Vzd�lenost bodu $x$ od mno�iny $A$} je ��slo $ d(x,A)=\inf\{ d(x,y)|y\in A\}$.
\end{definiceN}

\begin{definiceN}{uz�v�r}
\emph{Uz�v�rem} mno�iny $A$ v metrick�m prostoru $(M,d)$ naz�v�me mno�inu
$$\overline{A} = \bigcap_{\forall F}\{F\;\vert\;F \text{ uzav�en�}, A \subseteq F \subseteq M\}$$
Ekvivalentn� $\overline{A} = \lbrace a \in M\;\vert\;a = \lim_{n \to \infty} a_n$ pro n�jakou $(a_n) \subset A \rbrace$, neboli $\overline{A} = A \cup \lbrace \text{limitn� body }A\rbrace$.
\end{definiceN}

\begin{vetaN}{vlastnosti uz�v�ru}
Bu� $(M,d)$ metrick� prostor a $A,B$ mno�iny v n�m, potom plat�:
\begin{penumerate}
\item $\overline{\emptyset} = \emptyset$, $\overline{M} = M$
\item $A \subset B \Rightarrow \overline{A} \subset \overline{B}$
\item $\overline{\overline{A}} = \overline{A}$
\item $\overline{A \cup B} = \overline{A} \cup \overline{B}$
\item charakteristika uz�v�ru: $\overline{A} = \{x \in M| d(x,A) = 0\}$
\item $\overline{A}$ je nejmen�� uzav�en� mno�ina obsahuj�c� $A$, proto $A$ je uzav�en� pr�v� kdy� $\overline{A}=A$.
\end{penumerate}
\end{vetaN}

\begin{poznamkaN}{ekvivalentn� definice uzav�en� mno�iny}
Mno�ina $X\subset(M, d)$ je uzav�en�, jestli�e ka�d� posloupnost bod� $(x_n)_{n\geq 0}$, kter� v $X$ le�� a konverguje, v $X$ m� tak� svou limitu.
\end{poznamkaN}

\subsubsection*{Spojitost a stejnom�rn� spojitost}

\begin{definiceN}{spojit� zobrazen�}
Pro metrick� prostory $(X, d_1)$ a $(Y, d_2)$ je zobrazen� $f:X\to Y$ \emph{spojit� v bod�} $x\in X$, jestli�e
$$\forall\varepsilon>0\ \exists\delta>0:  d_1(x,y)<\delta\;\Rightarrow\;  d_2(f(x),f(y))<\varepsilon$$
Ekvivalentn� pro $\forall$ okol� $V$ bodu $f(x)$ $\exists$ okol� $U$ bodu $x$ tak, �e $y \in U \Rightarrow f(y) \in V$.\\
Zobrazen� $f$ je \emph{spojit�}, pokud je spojit� v ka�d�m bod� $x\in X$. 
\end{definiceN}

\begin{vetaN}{vlastnosti spojitosti}
Nech� je d�no zobrazen� $f:X\to Y$ mezi dv�ma metrick�mi prostory $(X, d_1)$, $(Y, d_2)$. Pak jsou n�sleduj�c� tvrzen� ekvivalentn�.
\begin{penumerate}
    \item $f$ je spojit�.
    \item Pro ka�dou konvergentn� posloupnost $(x_n)_{n\geq 0}$ v $X$ plat� $f(\lim x_n)=\lim f(x_n)$.
    \item Pro ka�d� $x$ a ka�d� okol� $U$ bodu $f(x)$ existuje okol� $V$ bodu $x$ takov�, �e $f[V]\subseteq U$.
    \item Je-li $U \subset Y$ otev�en� v $Y$, pak $f^{-1}(U) \subset X$ je otev�en� v $X$.
    \item Je-li $U \subset Y$ uzav�en� v $Y$, pak $f^{-1}(U) \subset X$ je uzav�en� v $X$.
    \item Pro ka�dou $A\subseteq X$ plat� $f[\overline{A}]\subseteq \overline{f[A]}$.
\end{penumerate}
\end{vetaN}

\begin{definiceN}{stejnom�rn� spojit� zobrazen�}
�ekneme, �e zobrazen� $f:(X, d_1)\to(Y, d_2)$ je \emph{stejnom�rn� spojit�}, jestli�e
$$\forall\varepsilon>0\ \exists\delta>0\text{ takov�, �e }\forall(x,y)\text{ plat� } d_1(x,y)<\delta\;\Rightarrow\; d_2(f(x),f(y))<\varepsilon$$
\end{definiceN}

\begin{vetaN}{skl�d�n� zobrazen�}
Slo�en� dvou spojit�ch (nebo stejnom�rn� spojit�ch) zobrazen� je spojit� (resp. stejnom�rn� spojit�).
\end{vetaN}

\begin{definiceN}{homeomorfismus}
Zobrazen� $f:(X,d_1)\to(Y,d_2)$ nazveme \emph{(stejnom�rn�m) homeomorfismem}, pokud
\begin{pitemize}
    \item $f$ je bijekce,
    \item $f$ je (stejnom�rn�) spojit�,
    \item inverzn� zobrazen� $f^{-1}$ je (stejnom�rn�) spojit�.
\end{pitemize}
Prostory $X$ a $Y$ pak jsou \emph{(stejnom�rn�) homeomorfn�}. Pokud je $f$ identick� zobrazen� z $(X, d_1)$ do $(X, d_2)$, ��k�me, �e metriky $d_1$ a $ d_2$ jsou \emph{(stejnom�rn�) ekvivalentn�}. (Jin� definice ekvivalentn�ch metrik po�aduje, aby metrick� prostory $(X, d_1)$ a $(X, d_2)$ m�li tyt� otev�en� mno�iny).
\end{definiceN}

\begin{vetaN}{aritmetika zobrazen�}
Jsou-li $f,g$ spojit� funkce $(X, d)\to\mathbb{R}$ (kde $(X, d)$ je metrick� prostor) a $\alpha\in\mathbb{R}$, potom i funkce $f+g$, $\alpha\cdot f$, $f\cdot g$ a $\frac{f}{g}$ (m�-li tato smysl) jsou spojit�. Plat� i pro spojitost zobrazen� v n�jak�m bod� $x_0\in X$.

\begin{dukaz}
D�kaz t�to v�ty je vlastn� stejn� jako d�kaz v�ty o aritmetice limit pro re�ln� funkce (jen pracujeme se zobrazen�mi na metrick�ch prostorech).
\end{dukaz}
\end{vetaN}

\subsubsection*{Podprostor metrick�ho prostoru}

\begin{definiceN}{podprostor}
Pro metrick� prostor $(X, d)$ a mno�inu $X_1\subset X$ vezmeme funkci $d_1:X_1\times X_1\to\mathbb{R}$ danou p�edpisem $d_1(x,y)= d(x,y)\ \forall x,y\in X_1$. Pak $(X_1, d_1)$ je \emph{podprostor} metrick�ho prostoru $(X, d)$ (\emph{indukovan�} podmno�inou $X_1$).
\end{definiceN}

\begin{poznamka}
Pro $X_1 \subset X$ je zobrazen� vlo�en� $j:(X_1, d)\to(X, d)$, definovan� p�edpisem $j(x)=x$ stejnom�rn� spojit�.
\end{poznamka}

\begin{vetaN}{vlastnosti podprostor�}
Bu� $Y$ podprostor metrick�ho prostoru $X$. Potom plat�
\begin{penumerate}
    \item $B_Y(x,\varepsilon)=B_X(x,\varepsilon)\cap Y$,\hfill\emph{(o okol� bod�)}
    \item $U$ je otev�en� v $Y$, pr�v� kdy� existuje otev�en� $V \subset X$ takov�, �e $U=V\cap Y$ (to sam� plat� i pro uzav�en� mno�iny),
    \item $\overline{A}^Y=\overline{A}^X\cap Y$.\hfill\emph{(o uz�v�ru mno�iny)}
\end{penumerate}
\end{vetaN}

\begin{vetaN}{podprostor zachov�v� spojitost}
Pro $f:(X, d_1)\to (Y, d_2)$ (stejnom�rn�) spojit� zobrazen� a $X_1\subseteq X$, $Y_1\subseteq Y$ takov�, �e $f[X_1]\subseteq Y_1$ je $f_1:X_1\to Y_1$ definovan� p�edpisem $f_1(x)=f(x)$ (stejnom�rn�) spojit�.
\end{vetaN}
\newpage

\subsection{Kompaktn� metrick� prostory}

\begin{definiceN}{kompaktn� prostor}
�ekneme, �e metrick� prostor je \emph{kompaktn�}, jestli�e v n�m lze z ka�d� posloupnosti bod� vybrat konvergentn� podposloupnost.
\end{definiceN}

\begin{priklady}
\begin{pitemize}
    \item Ka�d� kone�n� prostor je kompaktn�.
    \item Ka�d� omezen� uzav�en� interval (v $\mathbb{R}$!) je kompaktn�.
\end{pitemize}
\end{priklady}

\begin{vetaN}{zachov�n� kompaktnosti}
Kompaktnost se zachov�v�:
\begin{penumerate}
    \item p�echodem k uzav�en�mu podprostoru,
    \item obrazem spojit�m zobrazen�m,
    \item kart�zsk�m sou�inem.
\end{penumerate}

\begin{ideadukazu}
    1. -- m�m vybranou s~limitou, mus� z�st�vat v�echny limity v n�m (d�ky uzav�enosti). 2. -- dok�u za chv�li v~jin� v�t� :-) 
    
    3. -- u kone�n� mnoha bych vybral z~ka�d� konvergentn�, aby byly se stejn�mi indexy (\uv{postupn�}), a jejich sou�in by m�l b�t asi taky konvergentn�. U nekone�n�ho sou�inu nev�m.
\end{ideadukazu}

\end{vetaN}

\begin{definiceN}{omezen� podmno�ina}
Podmno�ina $A$ metrick�ho prostoru $(M,d)$ je \emph{omezen�}, pokud existuje kone�n� $K$ takov�, �e 
$$x,y\in A \;\Rightarrow\; d(x,y)<K$$
\end{definiceN}

\begin{vetaN}{uzav�enost kompaktn�ho prostoru}
Ka�d� kompaktn� podprostor $Y$ libovoln�ho metrick�ho prostoru $X$ je uzav�en� a omezen� (v $X$).

\begin{ideadukazu}
    \begin {itemize}
\item Uzav�enost: Z definice (kompaktn� prostory v�dy maj� konvergentn� podposloupnosti a tud� v nich le�� i p��slu�n� limity). 
\item Omezenost: sporem, nech� nen� omezen�. 

Pokud nen� omezen�, m��eme pro n�jak� pevn� dan� bod $a$ pro libovoln� velk� $K$ vz�t bod $b\in X$, �e $d(a,b)>K$ -- vezmeme dva body, vzd�len� od sebe $3K$, a d�ky troj�heln�kov� nerovnosti alespo� jeden z~nich je od $a$ vzd�len $K$.

Pak vytvo��me induktivn� posloupnost $(a_n)$ bez konvergentn� podposloupnosti. Za�n�me v~libovoln�m $a_1$. M�me $k$ prvn�ch �len�, vezmu $R=\max\limits_{i \in \left\lbrace 2 .. k \right\rbrace}\left\lbrace d(a_1, a_i)\right\rbrace + 1 $. Jak jsem uvedl v��e, existuje $a_{i+1}$, �e $d(a,b)>R$. D�ky troj�heln�kov� nerovnosti pak $\forall i \neq j:\;d(a_i, a_j) \geq 1$ a proto $(a_n)$ nem� konvergentn� podposloupnost.
    \end{itemize}
\end{ideadukazu}
\end{vetaN}

\begin{vetaN}{spojit� zobrazen� a kompaktn� mno�iny}
Bu� $f:(X,d_1)\to(Y,d_2)$ spojit� zobrazen� a $X$ kompaktn� metrick� prostor. Potom plat�:
\begin{penumerate}
    \item Zobrazen� $f$ je stejnom�rn� spojit�.
    \item $f[X]$ je kompaktn� podmno�ina $Y$.
    \item Je-li $f$ nav�c bijekce, pak je inverzn� zobrazen� $f^{-1}$ nutn� spojit�.

    \item Je-li $f$ nav�c bijekce, je $f$ stejnom�rn� homeomorfismus.

\end{penumerate}

\begin{ideadukazu}
    \begin {enumerate}
        \item Sporem. Kdyby nebylo stejnom�rn� spojit�, existovalo by $\varepsilon$, �e pro ka�d� $\delta$ existuj� body $x, y$, kter� jsou si bl�, ne� $\delta$, ale jejich zobrazen� jsou od sebe minim�ln� $\varepsilon$ (negace definice).
        
        $\delta$ m��eme po��d zmen�ovat, dostaneme tak dv� posloupnosti ($\left\{x_n\right\}$ a $\left\{y_n\right\}$), vybereme z~$\left\{x_n\right\}$ konvergentn� a z~$\left\{y_n\right\}$ tak�, tak, aby m�ly stejn� indexy $k_n$ (nejd��v vybereme z~jedn� konvergentn�, z~druh� vybereme tyt� indexy, z~n� je�t� znovu vybereme konvergentn�) -- to m��eme ud�lat, proto�e jsme v~kompaktn�m prostoru.
        
        Potom $\lim x_{k_n} = \lim y_{k_n}$, ale $\lim f\left(x_{k_n}\right) \neq \lim f\left(y_{k_n}\right),$ proto�e jsou od sebe v�dy zobrazen� vzd�lena minim�ln� $\varepsilon$. To je ale spor s~jedn�m z~n�sledk� spojitosti zobrazen�.
        
        \item Vezm�me si posloupnost bod� v~$f[X]$, ty mus� m�t vzory v~$X$, z~t�ch vzor� lze vz�t vybranou konvergentn� (d�ky spojitosti), mus� b�t konvergentn� i limita jej�ho obrazu (z~n�sledk� spojitosti zobrazen�).
        
        \item D�vejme se na inverzn� zobrazen� $f^{-1}$. Chceme uk�zat, �e obraz uzav�en� mno�iny v~$X$ zobrazen�m ${f^{-1}}^{-1}$ je uzav�en�, z~�eho� by plynulo, �e zobrazen� $f^{-1}$ je prost� (d�ky jedn� z p�edchoz�ch v�t). 
        
        Vezm�me $M$ uzav�en� v~$X$. Je kompaktn� (proto�e ka�d� uzav�en� mno�ina $X$ je kompaktn�). Podle 2. je i $f[X]$ kompaktn�. Proto�e $f$ je bijekce, ${f^{-1}}^{-1}[X]=f[X]$, tedy je kompaktn� a (d�ky p�edchoz� v�t�) je omezen�.
        
        \item je tot�, co 3. (viz definice homeomorfismu).
        
    \end{enumerate}
\end{ideadukazu}

\end{vetaN}

\begin{vetaN}{omezen� euklidovsk� podprostor}
\begin{penumerate}
    \item Podprostor $X$ euklidovsk�ho prostoru dimenze $n$ ($\mathbb{E}_n$) je kompaktn�, pr�v� kdy� je uzav�en� a omezen�. 
    \item Kompaktn� podprostor $X\subseteq\mathbb{R}$ ($\equiv\mathbb{E}_1$) m� nejv�t�� a nejmen�� prvek.
\end{penumerate}
\end{vetaN}

\begin{dusledek}
Z bod� 2. p�edchoz�ch dvou tvrzen� plyne, �e spojit� re�ln� funkce nab�v� na kompaktn�m prostoru minima i maxima. (Co� plat� i pro spojit� re�ln� funkce na kone�n�m defini�n�m oboru, jeliko� kone�nost $\Rightarrow$ kompaktnost.)
\end{dusledek}
\newpage

\subsection{�pln� prostory}

\begin{definiceN}{cauchyovsk� posloupnost bod�}
Posloupnost $(x_n)$ bod� z metrick�ho prostoru $(M,d)$ budi� \emph{cauchyovsk�}, jestli�e
$$\forall\varepsilon>0\ \exists n_0: m,n\geq n_0 d(x_m,x_n)<\varepsilon$$
\end{definiceN}

\begin{poznamka}
Je-li posloupnost $(x_n)$ konvergentn�, pak je cauchyovsk�. Obr�cen� implikace obecn� neplat�.
\end{poznamka}

\begin{definiceN}{�pln� prostor}
Prostor je \emph{�pln�}, pokud v n�m ka�d� cauchyovsk� posloupnost konverguje.
\end{definiceN}

\begin{vetaN}{o podposloupnosti}
Pokud m� cauchyovsk� posloupnost n�jakou konvergentn� podposloupnost, pak konverguje sama.
\end{vetaN}

\begin{priklady}
\begin{pitemize}
    \item $\mathbb{R}$ je �pln� prostor (d�ky Bolzanov�-Cauchyov� podm�nce).
    \item Ka�d� kompaktn� prostor je �pln� (podle p�edchoz� v�ty).
    \item $\mathbb{E}_n$ je �pln� prostor (bez d�kazu; vy�aduje sou�iny prostor�).
    \item Prostor $C[a, b]$ funkc� spojit�ch na $[a, b]$ s maximovou metrikou je �pln�. Je-li toti� posloupnost $(f_n)$ cauchyovsk�, spl�uje stejnom�rnou Bolzanovu-Cauchyovu podm�nku a tedy na $[a, b]$ konverguje stejnom�rn� k jist� funkci $f$. Funkce $f$ je na $M$ spojit�, proto�e je stejnom�rnou limitou spojit�ch funkc�. Tedy $f \in \mathcal{C}(M)$ a v supremov� metrice m�me $\lim_{n \to \infty} f_n = f$.
    \item Op�t pro $\mathcal{C}[a, b]$ te� ov�em s integr�ln� metrikou. Vznikl� metrick� prostor nen� �pln�. Sestroj�me cauchyovskou posloupnost, kter� nem� limitu. Polo��me $a = -1, b = 1$ a uv��me funkce $$f_n (x) =
    \begin{cases}
    -1 & \text{ pro } -1 \leq x \leq -n^{-1}\\
    nx & \text{ pro } -n^{-1} \leq x \leq n^{-1}\\
    1  & \text{ pro } n^{-1} \leq x \leq 1.
    \end{cases}$$
    Pak $(f_n) \subset \mathcal{C}[-1, 1]$ a $(f_n)$ je cauchyovsk�, proto�e pro $m \leq n$ m�me $d(f_m, f_n) = \int_{-1}^{1} |f_m(x) - f_n(x)| dx \leq \int_{-1/m}^{1/m} 1 dx = 2/m$. Neexistuje v�ak funkce $f \in \mathcal{C}[-1, 1]$, pro n� by $f_n \to f$  pro $n \to \infty$. Takov� funkce $f$ by podle definice $f_n$ musela b�t na intervalu $[-1, 0)$ identicky rovna $-1$ a na intervalu $(0, 1]$ identicky rovna $1$, co� je pro funkci spojitou na $[-1, 1]$ nemo�n�.
    \item Uva�me euklidovsk� metrick� prostory $\mathbb{R}$ a $(-\pi /2, \pi /2)$. Bijekce $f(x) = \arctan(x): \mathbb{R} \to (-\pi /2, \pi /2)$ je homeomorfismus, $f$ i $f^{-1}(x) = \tan(x): (-\pi /2, \pi /2) \to \mathbb{R}$ jsou spojit� zobrazen�. Ov�em $\mathbb{R}$ je �pln� metrick� prostor, ale $(-\pi /2, \pi /2)$ nikoli. �plnost metrick�ho prostoru nen�, na rozd�l od kompaktnosti, topologick� vlastnost, nen� ur�ena pouze otev�en�mi mno�inami, z�vis� i na metrice. Nicm�n� se �plnost se zachov�v� homeomorfismem, kter� je v obou sm�rech stejnom�rn� spojit� (funkce $\tan x$ nen� na $(-\pi /2, \pi /2)$ stejnom�rn� spojit�).
\end{pitemize}
\end{priklady}

\begin{vetaN}{zachov�n� �plnosti}
�plnost metrick�ho prostoru se zachov�v�:
\begin{penumerate}
    \item obrazem stejnom�rn� spojit�m prost�m zobrazen�m, pokud je i inverzn� zobrazen� stejnom�rn� spojit�,
    \item p�echodem k uzav�en�mu podprostoru,
    \item kart�zsk�m sou�inem.
\end{penumerate}

\begin{poznamka}
Pou��v� se zejm�na p�i nahrazov�n� metriky metrikou s n� stejnom�rn� ekvivalentn�. Tvrzen� pro \uv{oby�ejnou} spojitost neplat�.
\end{poznamka}
\end{vetaN}

\begin{vetaN}{o �pln�m podprostoru}
Podprostor $Y$ �pln�ho prostoru $X$ je �pln�, pr�v� kdy� je $Y$ v $X$ uzav�en� mno�ina.
\end{vetaN}

\begin{dusledek}
Kompaktn� metrick� prostor je v�dy �pln�.
\end{dusledek}


\end{document}
